\section{Introduction}
\label{sec:introduction}

Many scientific machine learning problems have objectives with \(m\) prescribed
components,
\begin{equation}
F(\theta)=\sum_{i=1}^{m}E_i(\theta).
\label{eq:intro-objective}
\end{equation}
For example, physics-informed neural network (PINN) objectives may combine
residual, boundary, initial-condition, and data terms
\cite{raissi2019pinn,karniadakis2021physics}, with regularization added when
needed. In operator learning such as DeepONet \cite{lu2021deeponet} , a data term can likewise be separated from
weight decay. These components can have different scales and gradient
directions. Once their gradients are summed, part of this directional
information may be hidden by aggregation or cancellation. Imbalance between
the residual and boundary terms of a PINN objective, and its effect on
training, motivate the adaptive loss weighting developed for physics-informed
models \cite{wang2021gradientpathologies,wang2022pinnntk,mcclenny2023selfadaptive}.
Those methods rescale the components before they are summed. The method
developed here keeps the sum in \cref{eq:intro-objective} as it stands and
uses the component gradients to build the curvature inside the solve.

Most first-order optimization methods act primarily on the aggregated
gradient. Momentum incorporates information from previous iterations
\cite{polyak1964}, while adaptive methods such as Adam and AMSGrad introduce
coordinatewise scaling \cite{kingma2015adam,reddi2018amsgrad}. Neither
mechanism directly uses the prescribed component structure of the objective.
In an energy-dissipative scheme, momentum adds cross-iteration terms and
adaptive scaling changes the quadratic form in which those terms are measured.

The scalar auxiliary variable (SAV) approach constructs energy-stable schemes
\cite{shen2018sav,shen2019savreview}. Originally developed for dissipative
gradient flows, SAV and its relaxed variants introduce auxiliary variables so
that a modified energy satisfies a discrete dissipation law
\cite{jiang2022rsav,zhang2022gsavrelax}. These ideas have been carried to
finite-dimensional optimization \cite{liu2022,liu2023savopt}, including a
relaxed vector auxiliary variable algorithm for unconstrained problems
\cite{zhang2025rvsavopt, ZHANGop1, ZHANGop2}. Standard SAV methods,
however, primarily depend on the total gradient and therefore do not fully
exploit the decomposition in \cref{eq:intro-objective}.

The pullback-corrected SAV (PB--SAV) family for gradient flows
\cite{zhang2026pbsav} separates energy tracking from the curvature
correction. A single scalar auxiliary variable tracks the total
shifted objective, while the component gradients generate the positive
semidefinite matrix
\[
B_m
=
\sum_{i=1}^{m}
\frac{\nabla E_i(\theta)\nabla E_i(\theta)^{\mathsf T}}
{2\bigl(E_i(\theta)+C_i\bigr)}.
\]
Here each shift \(C_i\) is chosen so that \(E_i(\theta)+C_i>0\). The matrix
\(B_m\) retains the component gradient directions and has rank at most \(m\).
Its Gauss--Newton interpretation and the corresponding approximation of
quadratic Hessians are established in \cite{zhang2026pbsav}. In a direct
Euclidean step, the aggregate SAV correction only rescales the gradient, while the component pullback can also change its direction. The correction therefore uses only component energies and
first derivatives.
\Cref{sec:background} restates that construction. We build on it an
optimizer with momentum and adaptive scaling, and determine how the component
curvature enters its energy law and its local stability condition.

Our main contributions are as follows.
\begin{itemize}[leftmargin=*]

\item \textbf{Component curvature and local stability.}
\Cref{prop:block-compression} characterizes the Gauss--Newton curvature
retained by a vector-valued least-squares component. For fixed mobility,
\cref{thm:local-schur-stability,cor:local-step-interval} give the exact local
stability condition and quantify the change in admissible step sizes relative
to the aggregate correction, including the
case in which the component gradients cancel, so that the aggregate correction
vanishes while the component correction does not.

\item \textbf{Momentum inside the solve.}
The gradient and stored momentum enter one implicit equation.
\Cref{prop:inertial-model} identifies its quadratic model and the excess
incurred by appending momentum after the solve, and
\cref{prop:appended-momentum-defect} shows that the corresponding identity
for the appended construction carries a cross term of indefinite sign. For the momentum update,
\cref{thm:monotone-mobility-dissipation} establishes an exact modified energy
law for positive definite mobilities that are nonincreasing in the Loewner
order, including Euclidean and AMSGrad-type choices.

\item \textbf{Convergence of the relaxed update.}
Under the local stability condition,
\cref{thm:local-relaxed-convergence} proves local geometric convergence near a
nondegenerate minimum for fixed mobility and every scalar relaxation
sequence. For varying mobility,
\cref{thm:conditional-stationarity,cor:best-iterate-stationarity} give
stationarity and an \(O(N^{-1/2})\) best-iterate gradient bound under uniform
mobility and curvature bounds, a positive lower bound on \(q^n/Q^n\), and
relaxation parameters bounded away from one.

\end{itemize}

\Cref{sec:background} introduces SAV and the PB--SAV geometry.
\Cref{sec:methodology} presents the method and analysis; the implementation
in \cref{alg:pbsav} uses a dense solve of order at most \(m\)
(\cref{prop:low-rank-solution}). \Cref{sec:experiments} reports five numerical
experiments and \cref{sec:conclusion} concludes. Longer proofs
appear in \ref{app:convergence-proofs}. \ref{app:directional-coverage} gives a
curvature condition for exact scalar tracking under full relaxation.
